\chapter{Développeur quant}\label{ch:in:quant-developer}

Pour l'exercice 2025, les employeurs de la finance référencés ont déposé 169 demandes de certification des conditions de travail pour des postes dont l'intitulé désignait
un développeur ou un ingénieur quantitatif. Ils en ont déposé 122 sous le code de métier des analystes quantitatifs financiers,
25 sous celui des développeurs de logiciels, et 22 sous au moins trois autres codes, dont aucun n'atteint dix. Le code relève du choix
de l'employeur, et il suit davantage le poste que l'intitulé~: derrière «~quant developer~» se cachent trois métiers différents,
l'\omterm{def:in:quant-developer:research}{ingénieur de recherche}, le \omterm{def:in:quant-developer:library}{développeur de bibliothèque} du desk et le \omterm{def:in:quant-developer:trading}{développeur de systèmes de trading}. Ce chapitre les distingue,
montre ce que chacun livre et comment il est jugé, et mesure ce que le choix du code révèle sur la rémunération.

\begin{center}\footnotesize
\begin{tabular}{@{}p{2.4cm}p{3.3cm}p{3.3cm}p{3.3cm}@{}}
\toprule
\multicolumn{4}{@{}l}{\textbf{Fiches de rôle~: les trois développeurs quant}}\\
\midrule
 & \omterm{def:in:quant-developer:research}{ingénieur de recherche} & \omterm{def:in:quant-developer:library}{développeur de bibliothèque} & \omterm{def:in:quant-developer:trading}{développeur de systèmes de trading}\\
\midrule
travaille avec & chercheurs & \omterm{def:in:bank-quant:strat}{stratégistes de desk} et \omterm{def:in:bank-quant:library}{quants de bibliothèque} & traders et ingénieurs\\
livre & l'infrastructure de recherche et les versions de production des signaux & le code de la bibliothèque de pricing & le système de trading\\
horizon & de quelques jours à quelques semaines & versions publiées & de la microseconde au jour\\
rattaché à & responsable de la recherche & responsable de l'analytique quantitative & responsable de la technologie de trading\\
codes dans les dépôts & 13-2099.01, 15-1252 & 13-2099.01, 15-1252 & 15-1252\\
enseigné dans & livre~7, chap.~1~; livre~15, chap.~1 & livre~5, chap.~28~; livre~6, chap.~29 & livre~13, chap.~1~; livre~16, chap.~21\\
\bottomrule
\end{tabular}
\end{center}

Le livre~16, chapitre~21, définit le développeur quant comme un ingénieur qui transforme le prototype d'un chercheur en code de production,
il en répond pour la performance et la correction, connaît assez le modèle pour le modifier sans risque, et place le rôle dans des équipes intégrées
ou centrales. Ce chapitre retient cette définition et se demande lequel des trois produits l'ingénieur est en train de construire.

\section{L'ingénieur de recherche}

\begin{definition}[Ingénieur de recherche]\label{def:in:quant-developer:research}
Un \emph{ingénieur de recherche}\index{ingénieur de recherche} est un développeur quant qui travaille au sein d'un groupe de recherche et construit
ce dont se servent les chercheurs et ce que deviennent leurs résultats~: pipelines de données, outils de backtest et de simulation, le
calcul qui fait tourner les expériences, et les versions de production des signaux que le groupe a décidé de trader.
\end{definition}

Les clients de l'\omterm{def:in:quant-developer:research}{ingénieur de recherche} sont les chercheurs, et le produit est leur vitesse. Le livre~7 (chapitre~1) décrit le
processus de recherche que sert l'ingénieur, et le livre~15 les plateformes (recherche, données et risque) qu'une société construit pour lui~:
le magasin de données point-in-time, l'outil de backtest qui reproduit un résultat un an plus tard, l'outil de suivi d'expériences (experiment tracker), et
l'ordonnanceur de tâches sur le cluster. Dans un fonds systématique ou une plateforme, où de nombreux chercheurs partagent une même infrastructure, le
travail de l'\omterm{def:in:quant-developer:research}{ingénieur de recherche} est démultiplié par chaque chercheur qui l'utilise~; chez un teneur de marché, où la recherche est proche du
trading, l'ingénieur porte aussi un signal du code de recherche jusqu'au système de trading et vérifie que les deux concordent.

Le poste est jugé sur deux points. Le premier est la correction en reproduction~: un backtest qui ne peut pas être relancé, ou un
signal de production qui diffère de sa version de recherche détruit les preuves que la recherche a pour mission de produire. Le second
est le temps~: les heures qui séparent une idée de son premier test honnête. La \omterm{def:in:quant-researcher:feedback}{vitesse du retour d'information} du chercheur (chapitre~17) est limitée par
le marché~; l'\omterm{def:in:quant-developer:research}{ingénieur de recherche} fixe l'autre limite.

\section{Le développeur de bibliothèque du desk}

\begin{definition}[Développeur de bibliothèque]\label{def:in:quant-developer:library}
Un \emph{développeur de bibliothèque}\index{développeur de bibliothèque} est un développeur quant qui écrit et maintient le code
d'une bibliothèque de pricing et de risque destinée aux quants qui conçoivent ses modèles~: les interfaces, structures de données, performance et
l'outillage numérique qui transforme les méthodes d'un \omterm{def:in:bank-quant:library}{quant de bibliothèque} en un produit que chaque desk et chaque système de risque peut appeler.
\end{definition}

Le \omterm{def:in:quant-developer:library}{développeur de bibliothèque} travaille aux côtés du \omterm{def:in:bank-quant:library}{quant de bibliothèque} du chapitre~18, et la frontière entre les deux tient surtout à l'accent
mis~: le quant porte les mathématiques d'un modèle, le développeur le code qui le rend rapide, stable et appelable depuis
ailleurs. Le livre~5, chapitre~28, construit une bibliothèque de pricing avec un cœur en C++20 et Rust derrière une interface Python, et la
répartition du travail reprend celle des deux métiers~: la méthode numérique d'un côté~; la disposition en mémoire, le parallélisme et
l'interface, la compilation et la suite de non-régression de l'autre. Lorsqu'une seule bibliothèque sert chaque desk d'une banque, le poste
relève du groupe qui la possède~; lorsque le pricing fait partie du système de trading, comme pour un teneur de marché d'options (livre~11), il
se confond avec celui du \omterm{def:in:quant-developer:trading}{développeur de systèmes de trading}.

Le \omterm{def:in:quant-developer:library}{développeur de bibliothèque} est jugé comme la bibliothèque elle-même~: aucun chiffre d'un desk ne change sans que quelqu'un l'ait voulu, chaque
version reproduit les résultats de la précédente à une tolérance énoncée, et la bibliothèque est assez rapide pour les calculs de risque qui
l'appellent des millions de fois par nuit.

\section{Le développeur de systèmes de trading}

\begin{definition}[Développeur de systèmes de trading]\label{def:in:quant-developer:trading}
Un \emph{développeur de systèmes de trading}\index{développeur de systèmes de trading} est un développeur quant qui construit et fait tourner les systèmes
qui tradent~: gestionnaires de données de marché, moteur de stratégie, saisie des ordres, contrôles de risque et supervision qui les entoure, mesurés
en microsecondes et par leur comportement lorsqu'une panne survient.
\end{definition}

Le \omterm{def:in:quant-developer:trading}{développeur de systèmes de trading} travaille au plus près de l'argent. La latence tick-to-trade et le budget de latence du livre~13,
chapitre~1, sont les unités de ce développeur~; les protocoles des bourses du livre~10 et le réseau et le matériel du livre~14 sont le
cadre de travail. La part quantitative du poste n'est pas le modèle mais son implémentation sous contraintes~: un
signal qui doit être calculé en deux microsecondes, un biais d'inventaire qui doit être appliqué à chaque cotation, un contrôle de risque
qui doit rejeter un ordre avant qu'il ne parte. Chez les teneurs de marché et les acteurs du trading à haute fréquence, ces développeurs sont souvent classés comme
ingénieurs logiciels, et le chapitre~20 traite de l'ingénierie qui les entoure.

\begin{dated}{2025-09}{Développeurs quant dans les dépôts et l'enquête}\label{dat:in:quant-developer:codes}
Demandes de certification des conditions de travail pour l'exercice 2025 chez les employeurs de la finance référencés~: 169 dans la famille d'intitulés «~quant developer~» (122
codés 13-2099.01, 25 codés 15-1252, 22 sous des codes comptant chacun moins de dix)~; la cellule des fonds systématiques, 12 demandes
émanant de 3 employeurs, \omterm{def:in:how-pay-works:base}{salaire de base} médian proposé 235\,000~\$ (10\up{e}--90\up{e}~centile 195\,500~\$--272\,500)~; chaque autre type d'employeur
dont la cellule provient d'un ou deux employeurs et est supprimée. Tous intitulés confondus~: développeurs de logiciels (15-1252) 5\,531
demandes, salaire médian 155\,000~\$~; analystes quantitatifs financiers (13-2099.01) 1\,294, salaire médian 150\,000~\$. Enquête sur les
métiers, mai 2025, développeurs de logiciels~: secteur des valeurs mobilières, salaire médian 163\,290~\$ (10\up{e}--90\up{e}~centile 104\,600~\$--245\,380~; 34\,940
salariés)~; banques 136\,620~\$~; éditeurs de logiciels 164\,550~\$ (173\,080 salariés).
\end{dated}

\section{Ce que livre chacun et comment il est jugé}

Les trois métiers partagent un intitulé et une bonne part d'une boîte à outils (un langage système, Python, gestion de versions, tests, quelques
notions de statistiques), mais ni le même produit ni la même horloge.
\begin{center}\footnotesize
\begin{tabular}{@{}p{2.7cm}p{3.3cm}p{3.1cm}p{3.3cm}@{}}
\toprule
 & \omterm{def:in:quant-developer:research}{ingénieur de recherche} & \omterm{def:in:quant-developer:library}{développeur de bibliothèque} & \omterm{def:in:quant-developer:trading}{développeur de systèmes de trading}\\
\midrule
livre & pipelines, outils, signaux de production & versions publiées de la bibliothèque & versions publiées du système, configurations\\
jugé sur & reproductibilité~; temps des chercheurs & stabilité~; vitesse~; couverture & latence~; disponibilité~; comportement correct\\
la panne ressemble à & un résultat que personne ne peut relancer & les chiffres de tous les desks bougent & une perte en quelques secondes\\
astreinte & rarement & à chaque version & pendant que les marchés sont ouverts\\
\bottomrule
\end{tabular}
\end{center}

Les codes professionnels suivent la même logique. La définition des analystes financiers quantitatifs, «~Develop mathematical or
statistical models for risk management, asset optimization, pricing, or relative value analysis~», correspond à l'\omterm{def:in:quant-developer:research}{ingénieur de recherche} et
au \omterm{def:in:quant-developer:library}{développeur de bibliothèque}~; celle des développeurs de logiciels, «~Research, design, and develop computer and network
software~», correspond au \omterm{def:in:quant-developer:trading}{développeur de systèmes de trading}. O*NET place la première dans sa bande de préparation la plus élevée (la plupart exigent des
études supérieures) et la seconde dans la bande inférieure (la plupart exigent un diplôme de premier cycle). Le code évolue aussi~: en 2018, la classification
a fusionné les deux codes de développeurs de logiciels de 2010 (15-1132 et 15-1133) dans le code 15-1252, si bien que les dépôts de l'exercice fiscal 2021 ont besoin
d'une table de correspondance avant d'être comparés à ceux de l'exercice fiscal 2025.

\begin{omfigure}
\begin{tikzpicture}
\begin{axis}[xbar stacked,width=10.5cm,height=5.2cm,xmin=0,xmax=100,ymin=0.4,ymax=4.6,bar width=11pt,ytick=data,
  yticklabels={{ML et données},{ingénieur logiciel},{développeur quant},{chercheur quantitatif}},yticklabel style={font=\scriptsize},
  xticklabel style={font=\footnotesize},xlabel={\small part des dépôts de l'exercice fiscal 2025,~\%},table/col sep=comma,
  legend style={legend cell align=left,font=\scriptsize,at={(0.5,-0.32)},anchor=north,/tikz/every even column/.append style={column sep=6pt}},
  legend columns=5]
\addplot[fill=omThm,draw=none] table[x=quant,y=pos]{figdata/industry/19-quant-developer/soc_mix.csv};
\addplot[fill=omDef,draw=none] table[x=software,y=pos]{figdata/industry/19-quant-developer/soc_mix.csv};
\addplot[fill=omDef!40,draw=none] table[x=stats,y=pos]{figdata/industry/19-quant-developer/soc_mix.csv};
\addplot[fill=omThm!40,draw=none] table[x=data,y=pos]{figdata/industry/19-quant-developer/soc_mix.csv};
\addplot[fill=black!25,draw=none] table[x=other,y=pos]{figdata/industry/19-quant-developer/soc_mix.csv};
\legend{13-2099.01,15-1252,15-2041 ou 15-2031,15-2051,autre}
\end{axis}
\end{tikzpicture}

\omcaption{Codes professionnels au sein de quatre familles d'intitulés, dépôts de l'exercice fiscal 2025 auprès des employeurs financiers sourcés. Codes
comptant moins de dix dépôts dans une famille sont comptés comme «~autre~». Les intitulés de développeur quant sont déposés surtout comme
analystes quantitatifs, les intitulés d'ingénieur logiciel presque uniquement comme développeurs de logiciels. Données~: \texttt{data/industry/lca\_titles.csv}
et \texttt{lca\_ranges.csv}, par l'intermédiaire de \texttt{in\_quantdev.soc\_mix}.}\label{fig:in:quant-developer:mix}
\end{omfigure}

Le code compte-t-il pour la rémunération~? Tous intitulés confondus chez les employeurs sourcés pendant l'exercice fiscal 2025, les dépôts classés comme développeurs de logiciels
ont un \omterm{def:in:how-pay-works:base}{salaire de base} médian proposé supérieur de 5\,000~\$ à ceux classés comme analystes quantitatifs, avec un intervalle bootstrap de
$-\$372$ à 5\,688~\$ qui contient tout juste zéro. Au sein de chaque \omterm{def:in:pay-levels-by-role-firm-type-and-seniority:soc}{niveau de salaire}, le signe s'inverse~: la médiane des analystes est plus élevée
aux niveaux III et IV, de 11\,000~\$ et 27\,092~\$ (\cref{fig:in:quant-developer:levels,fig:in:quant-developer:gap}).
Les deux constats sont compatibles, car les codes sont utilisés à des niveaux différents~: la moitié des dépôts de développeurs de logiciels
sont au niveau IV, contre un sur sept pour les analystes (\cref{fig:in:quant-developer:levelmix}).

\begin{omfigure}
\begin{tikzpicture}
\begin{axis}[width=10.5cm,height=5.2cm,xmin=0.7,xmax=4.3,ymin=80,ymax=215,xtick={1,2,3,4},xticklabels={I,II,III,IV},
  xlabel={\small niveau de salaire},ylabel={\small salaire de base médian proposé, milliers de \$},tick label style={font=\footnotesize},
  grid=major,grid style={gray!20},table/col sep=comma,
  legend style={font=\scriptsize,at={(0.03,0.97)},anchor=north west,legend cell align=left}]
\addplot[omDef,thick,mark=*,error bars/.cd,y dir=both,y explicit]
  table[x expr=\thisrow{level}-0.05,y=dev,y error plus expr=\thisrow{dev_hi}-\thisrow{dev},
  y error minus expr=\thisrow{dev}-\thisrow{dev_lo}]{figdata/industry/19-quant-developer/levels.csv};
\addplot[omThm,thick,mark=square*,error bars/.cd,y dir=both,y explicit]
  table[x expr=\thisrow{level}+0.05,y=quant,y error plus expr=\thisrow{quant_hi}-\thisrow{quant},
  y error minus expr=\thisrow{quant}-\thisrow{quant_lo}]{figdata/industry/19-quant-developer/levels.csv};
\legend{développeurs de logiciels (15-1252),analystes quantitatifs (13-2099.01)}
\end{axis}
\end{tikzpicture}

\omcaption{\omterm{def:in:how-pay-works:base}{Salaire de base} médian proposé par \omterm{def:in:pay-levels-by-role-firm-type-and-seniority:soc}{niveau de salaire} pour les deux codes sous lesquels un développeur quant est déposé, tous intitulés confondus auprès des
employeurs financiers sourcés, exercice fiscal 2025, avec des intervalles bootstrap à 95~\% pour la médiane. Données~: \texttt{data/industry/lca\_soc.csv},
par l'intermédiaire de \texttt{in\_quantdev.soc\_cells}.}\label{fig:in:quant-developer:levels}
\end{omfigure}

\begin{omfigure}
\begin{tikzpicture}
\begin{axis}[width=10.5cm,height=5cm,xmin=0.5,xmax=5.5,ymin=-40,ymax=15,xtick={1,2,3,4,5},
  xticklabels={tous niveaux,I,II,III,IV},xlabel={\small niveau de salaire},ylabel={\small écart, milliers de \$},
  tick label style={font=\footnotesize},grid=major,grid style={gray!20},table/col sep=comma,
  legend style={font=\scriptsize,at={(0.97,0.97)},anchor=north east,legend cell align=left}]
\addplot[only marks,mark=o,black!60,thick,error bars/.cd,y dir=both,y explicit]
  table[x expr=\thisrow{pos}-0.1,y=d2021,y error plus expr=\thisrow{hi2021}-\thisrow{d2021},
  y error minus expr=\thisrow{d2021}-\thisrow{lo2021}]{figdata/industry/19-quant-developer/gap.csv};
\addplot[only marks,mark=*,omDef,thick,error bars/.cd,y dir=both,y explicit]
  table[x expr=\thisrow{pos}+0.1,y=d2025,y error plus expr=\thisrow{hi2025}-\thisrow{d2025},
  y error minus expr=\thisrow{d2025}-\thisrow{lo2025}]{figdata/industry/19-quant-developer/gap.csv};
\draw[black!50] (axis cs:0.5,0) -- (axis cs:5.5,0);
\legend{exercice fiscal 2021,exercice fiscal 2025}
\end{axis}
\end{tikzpicture}

\omcaption{\omterm{def:in:how-pay-works:base}{Salaire de base} médian proposé des dépôts de développeurs de logiciels moins celui des dépôts d'analystes quantitatifs, global et par
\omterm{def:in:pay-levels-by-role-firm-type-and-seniority:soc}{niveau de salaire}, avec des intervalles bootstrap à 95~\% (codes de l'exercice fiscal 2021 fusionnés par la table de correspondance de 2018). Au global, il est supérieur ou égal à zéro~; au sein de chaque niveau, il est négatif ou
indiscernable de zéro. Données~: \texttt{data/industry/lca\_soc\_gap.csv}, par l'intermédiaire de
\texttt{in\_quantdev.gaps}.}\label{fig:in:quant-developer:gap}
\end{omfigure}

\begin{omfigure}
\begin{tikzpicture}
\begin{axis}[ybar,area legend,width=10.5cm,height=4.8cm,xmin=0.5,xmax=4.5,ymin=0,ymax=60,bar width=14pt,xtick={1,2,3,4},
  xticklabels={I,II,III,IV},xlabel={\small niveau de salaire},ylabel={\small part des dépôts,~\%},
  tick label style={font=\footnotesize},grid=major,grid style={gray!20},table/col sep=comma,
  legend style={legend cell align=left,font=\scriptsize,at={(0.5,-0.42)},anchor=north,/tikz/every even column/.append style={column sep=8pt}},
  legend columns=2]
\addplot[fill=omDef,draw=none] table[x=level,y=dev]{figdata/industry/19-quant-developer/level_mix.csv};
\addplot[fill=omThm,draw=none] table[x=level,y=quant]{figdata/industry/19-quant-developer/level_mix.csv};
\legend{développeurs de logiciels (15-1252),analystes quantitatifs (13-2099.01)}
\end{axis}
\end{tikzpicture}

\omcaption{Répartition des dépôts des deux codes de l'exercice fiscal 2025 sur les quatre niveaux de salaire (dépôts avec un niveau indiqué). La moitié des
dépôts de développeurs de logiciels sont au niveau IV, ce qui tire vers le haut leur médiane globale. Données~: \texttt{in\_quantdev.level\_mix}.}\label{fig:in:quant-developer:levelmix}
\end{omfigure}

Cette comparaison montre ce que les codes ne disent pas~: qu'un développeur est mieux payé du seul fait d'être classé comme développeur de logiciels. Au
même \omterm{def:in:pay-levels-by-role-firm-type-and-seniority:soc}{niveau de salaire} déclaré, le code des analystes porte l'offre la plus élevée~; globalement, le code des logiciels regroupe davantage de dépôts de rang supérieur,
et la plupart viennent des équipes d'ingénierie des banques (4\,790 sur 5\,531, issus de huit banques). La rémunération d'un développeur quant dépend de l'employeur et
du niveau, comme le chapitre~14 l'a montré~; le code n'est qu'une description, et les deux écarts sont un avertissement contre la comparaison de médianes
portant sur des populations aux compositions différentes.

\section{Passer de l'un à l'autre}

Les trois métiers sont assez proches pour que l'on passe de l'un à l'autre, et assez éloignés pour que le passage ait un coût.
\begin{itemize}
\item \textbf{Du développement de systèmes de trading à l'ingénierie de recherche} rapproche le développeur de la recherche~: les compétences en systèmes
se transposent, et un développeur qui a fait tourner la version de production d'un signal sait où le code de recherche casse.
\item \textbf{De l'ingénierie de recherche à la recherche} est un changement de produit, des outils vers les résultats. Un \omterm{def:in:quant-developer:research}{ingénieur de recherche} qui écrit déjà le
code de production des signaux peut prendre un problème de recherche qui lui est propre~; les entreprises ne l'autorisent pas toutes de la même façon
(chapitre~28).
\item \textbf{Du développement de bibliothèque au desk} suit le parcours de la banque décrit au chapitre~18~: un \omterm{def:in:quant-developer:library}{développeur de bibliothèque} qui connaît les modèles
devient \omterm{def:in:bank-quant:strat}{stratégiste de desk}~; celui qui préfère les systèmes passe à la technologie de trading.
\item \textbf{D'un employeur à l'autre} le code et l'intitulé se transposent mal~: un «~développeur quant~» dans une banque et chez un teneur
de marché peut assurer respectivement le travail de bibliothèque et celui de système de trading. Un candidat doit demander lequel des trois produits
le poste livre et en quoi consiste l'astreinte.
\end{itemize}

\section{Tutoriel~: un titre, plusieurs codes}\label{tut:in:quant-developer:tut}

\begin{tutorial}
\textbf{Objectif.} Séparer l'intitulé du poste et du code, et mesurer ce que le code dit de la rémunération.
\textbf{État final~:} \cref{fig:in:quant-developer:mix,fig:in:quant-developer:levels,fig:in:quant-developer:gap,fig:in:quant-developer:levelmix}.
\begin{tutsteps}
\item \textbf{Normaliser les intitulés.} \texttt{firm.roles.normalise\_title} retire les équipes et les lieux, enlève les mots de séniorité,
développe les abréviations et supprime les chiffres de grade (\cref{lst:in:quant-developer:normalise})~; sur la liste illustrative du chapitre,
«~Sr.\ Quant Developer - Equities~», «~VP, Quant Dev~» et «~Quant Dev L3~» deviennent
tous \texttt{QUANTITATIVE DEVELOPER}.
\omcode{code/firm/roles/firm_roles.py}{224}{243}{Le normaliseur d'intitulés de poste.}\label{lst:in:quant-developer:normalise}
\item \textbf{Croiser les tableaux.} \texttt{crosstab(pairs)} compte les familles d'intitulés par code de profession et supprime les cases
de moins de dix~; le fichier \texttt{lca\_titles.csv} du chapitre~14 est ce tableau pour les dossiers déposés.
\item \textbf{Effectifs par code et par niveau.} \texttt{in\_lca\_soc\_derive.py} reconstruit les tableaux par code à partir des enregistrements
mis en cache du chapitre~14 (sans second passage sur les fichiers bruts), avec la règle des trois employeurs.
\item \textbf{L'écart.} \texttt{median\_gap(x, y, rng)} donne la différence et le rapport des médianes, à partir d'intervalles de
bootstrap indépendants (\cref{lst:in:quant-developer:gap}).
\omcode{code/firm/roles/firm_roles.py}{254}{265}{L'écart entre deux médianes avec intervalles de bootstrap.}\label{lst:in:quant-developer:gap}
\end{tutsteps}
\end{tutorial}

Pour l'exercice fiscal 2025~: écart global de 5\,000~\$ ($-\$372$ à 5\,688~\$)~; niveau III $-\$11\,000$ ($-\$19\,038$ à $-\$6\,600$)~; niveau
IV $-\$27\,092$ ($-\$37\,500$ à $-\$19\,900$). À répartition des niveaux inchangée, la médiane par niveau des développeurs de logiciels
est en moyenne inférieure de 7\,789~\$ à celle des analystes.

\textbf{Pour aller plus loin.} Tenir compte de l'employeur (un écart calculé au sein de chaque employeur, puis moyenné)~; séparer
les nouvelles embauches à temps plein des mutations~; relier directement les intitulés normalisés aux codes, ce qui suppose les intitulés que
le passage unique sur les fichiers de ce chapitre n'a pas conservés.

\section{Construction~: normaliseur d'intitulés et comparaison des codes}\label{bld:in:quant-developer:roles}

\begin{build}
\bfield{Rôle} Ajouter les trois développeurs quant au registre des postes, ainsi que les outils pour étudier les intitulés et les codes.
\bfield{Interface} \texttt{firm.roles}~: les fiches \texttt{research engineer}, \texttt{library developer},
\texttt{trading-system developer}~; \texttt{normalise\_title(title) -> (clean, seniority)}~;
\texttt{crosstab(pairs, min\_cell)}~; \texttt{median\_gap(x, y, rng, n\_boot, min\_cell)}.
\bfield{Règles} La séniorité est enregistrée, sans être conservée dans l'intitulé~; la première partie d'un intitulé qui désigne un métier l'emporte~;
les cases de moins de dix sont supprimées~; un écart n'est pas calculé à partir de moins de dix valeurs de part et d'autre.
\bfield{Tests d'acceptation} \texttt{code/firm/roles/tests/}~: six intitulés normalisés à leurs formes attendues~; un croisement supprime une case
de trois~; un écart implanté de 50 se trouve dans son intervalle~; une population de neuf ne donne pas d'écart.
\bfield{Pour aller plus loin} Apprendre la table d'abréviations du normaliseur à partir des données~; stratifier l'écart par employeur~; un test de
permutation de l'écart.
\end{build}

\begin{omsources}
\item O*NET OnLine, professions 15-1252, 13-2099.01, 15-2041 et 15-2051 (CC BY 4.0).
\item Bureau of Labor Statistics, SOC 2010 to 2018 crosswalk~; enquête sur les professions, mai 2025.
\item Tableaux du chapitre~14 issus des fichiers LCA du Department of Labor.
\item Livre~16, chapitre~21, pour le développeur quant et les équipes intégrées et centrales.
\end{omsources}

\section{Exercices}

\begin{exercise}[$\star$]\label{exo:in:quant-developer:1}
Quelle part des dossiers de développeur quant de l'exercice fiscal 2025 a été codée comme analystes financiers quantitatifs, et quelle part comme
développeurs de logiciels~?
\end{exercise}

\begin{exercise}[$\star$]\label{exo:in:quant-developer:2}
Que renvoie \texttt{normalise\_title} pour «~Sr.\ Quant Dev (Rates), NY~»~?
\end{exercise}

\begin{exercise}[$\star$]\label{exo:in:quant-developer:3}
D'après l'encadré daté, de combien la médiane de l'enquête pour les développeurs de logiciels est-elle supérieure dans le secteur des valeurs mobilières
à celle des banques, exprimée en rapport~?
\end{exercise}

\begin{exercise}[$\star\star$]\label{exo:in:quant-developer:4}
Expliquer en trois phrases comment la médiane des développeurs de logiciels peut dépasser la médiane globale des analystes
tout en lui étant inférieure aux niveaux III et IV.
\end{exercise}

\begin{exercise}[$\star\star$]\label{exo:in:quant-developer:5}
Pondérer la médiane de chaque code par niveau à l'aide des parts de niveaux des analystes (7.8~\%, 51.1~\%, 26.7~\%, 14.4~\%). Quelle est
la différence~?
\end{exercise}

\begin{exercise}[$\star\star$]\label{exo:in:quant-developer:6}
Une diplômée veut devenir chercheur dans les trois ans. Lequel des trois postes de développeur quant offre le meilleur
départ, et que doit-elle demander lors de l'entretien~?
\end{exercise}

\begin{exercise}[$\star\star\star$]\label{exo:in:quant-developer:7}
\emph{Programmation.} Écrire une fonction qui indique qu'un écart est distinguable de zéro lorsque son intervalle ne contient pas zéro, et
l'appliquer aux dix écarts de \texttt{lca\_soc\_gap.csv}. Combien sont distinguables~?
\end{exercise}

\begin{exercise}[$\star\star\star$]\label{exo:in:quant-developer:8}
\emph{Trouver la faille.} «~Les développeurs de logiciels dans les employeurs financiers gagnent plus que les analystes quantitatifs~: leur médiane est
supérieure de 5\,000~\$. Un développeur quant devrait demander à être classé comme développeur de logiciels.~»
\end{exercise}

\section{Problème~: un titre, trois postes}

\begin{problem}[{Problème du week-end --- un titre, trois postes}]\label{pb:in:quant-developer:1}
Un développeur a trois offres, chacune intitulée «~développeur quantitatif~»~: dans le groupe d'analyse d'une banque, dans le groupe de recherche d'un fonds
systématique et dans le groupe de technologie de trading d'un teneur de marché. Il veut savoir ce que cache l'intitulé.

\medskip\noindent\textbf{Partie I --- Les postes.}
\begin{enumerate}
\item Définir le développeur quant (livre~16) et les trois postes.
\item Que livre chacun~?
\item Sur quoi chacun est-il évalué, et à quoi ressemble l'échec~?
\item Quel poste comporte une astreinte pendant que les marchés sont ouverts~?
\item Où se situe chacun dans l'organisation~?
\end{enumerate}

\medskip\noindent\textbf{Partie II --- Les codes.}
\begin{enumerate}[resume]
\item Que désignent les codes de profession 15-1252 et 13-2099.01, et quels postes correspondent à chacun~?
\item Quelle préparation O*NET associe-t-il à chacun~?
\item Qu'est-ce qui a changé dans les codes des développeurs de logiciels en 2018~?
\item Donner la répartition SOC des dossiers de développeur quant de l'exercice fiscal 2025.
\item Pourquoi les cases de rémunération de la plupart des types d'employeurs sont-elles supprimées pour cette famille~?
\end{enumerate}

\medskip\noindent\textbf{Partie III --- L'écart.}
\begin{enumerate}[resume]
\item Énoncer la méthode de \texttt{median\_gap}.
\item Donner l'écart global de l'exercice fiscal 2025 et son intervalle.
\item Donner les écarts aux niveaux III et IV.
\item Expliquer l'inversion.
\item Que devient l'écart lorsque la répartition par niveau est fixée à celle des analystes~?
\end{enumerate}

\medskip\noindent\textbf{Partie IV --- Le verdict.}
\begin{enumerate}[resume]
\item Énoncer le \emph{résultat nommé}~: la répartition SOC des dossiers de développeur quant et l'écart de \omterm{def:in:how-pay-works:base}{salaire de base} médian proposé
entre les codes des développeurs de logiciels et des analystes quantitatifs, avec son intervalle.
\item Le code sous lequel il est classé compte-t-il pour sa rémunération~?
\item Que montre la case des fonds systématiques, et pourquoi est-elle la seule~?
\item Quelle question sépare le mieux les trois offres~?
\item En deux phrases, lui répondre.
\end{enumerate}
\end{problem}

\section{Questions d'entretien}

\begin{interviewq}[$\star$ \iqroles{développeur}]\label{iq:in:quant-developer:1}
Le signal Python d'un chercheur et votre version C++ divergent à la quatrième décimale sur 0.1~\% des jours. Que faites-vous~?
\end{interviewq}

\begin{interviewq}[$\star$ \iqroles{développeur}]\label{iq:in:quant-developer:2}
Comment rendriez-vous un backtest reproductible un an plus tard~?
\end{interviewq}

\begin{interviewq}[$\star\star$ \iqroles{développeur, banque}]\label{iq:in:quant-developer:3}
Une fonction de pricing est appelée dix millions de fois lors du calcul de risque de nuit et prend 40 microsecondes. Où regardez-vous en premier~?
\end{interviewq}

\begin{interviewq}[$\star\star$ \iqroles{développeur, trader}]\label{iq:in:quant-developer:4}
La latence tick-to-trade médiane du système de trading est correcte, mais le 99.9\up{e}~centile a doublé après une mise en production. Comment trouvez-vous
la cause~?
\end{interviewq}

\begin{interviewq}[$\star\star$ \iqroles{développeur, chercheur}]\label{iq:in:quant-developer:5}
Que mettriez-vous dans l'interface d'un signal pour que la recherche et la production ne puissent pas diverger silencieusement~?
\end{interviewq}

\begin{interviewq}[$\star\star\star$ \iqroles{développeur}]\label{iq:in:quant-developer:6}
Concevez la suite de tests de régression d'une bibliothèque de pricing utilisée par de nombreux desks.
\end{interviewq}
